\section{Matrix variation along the filter}
\label{sec:matrix}

The proof of Proposition~\ref{prop:filtered} must control all the states
of a filtered Gaussian vector, even when the sequence is arbitrarily
long.  We first prove a deterministic estimate that measures variation
of the filter by a total energy, rather than by the number of steps.
Spectral localization will give separate bounds for changes in the
sequence index and in Ornstein--Uhlenbeck time.

Throughout this section, fix a recursion as in
Proposition~\ref{prop:filtered} and a positive definite diagonal matrix
$Q\in\R^{m\times m}$.  The sequence index is $t\in\{0,\ldots,n\}$.
For a real matrix $A$, write
$\lVert A\rVert_{\HS}^2=\tr(AA^\top)$.

\subsection{Trace budgets and spectral scales}

The support of $R_{t-1}$ is contained in the initial block and the first
$t-1$ new coordinates.  Hence the new coordinate has no cross term in
$R_tR_t^\top$, and induction gives
\begin{equation}
 R_tR_t^\top=I_d,
 \qquad
 L_t:=R_tQR_t^\top
   =C_tL_{t-1}C_t^\top+Q_{d+t,d+t}b_tb_t^\top.
 \label{eq:speed-recursion}
\end{equation}
In particular $L_t$ is positive definite.  Every $C_t$ is a contraction,
since $C_tC_t^\top\preceq I_d$.

Let $\ell_t$ be the trace lost in the contraction in
\eqref{eq:speed-recursion}, and let $g_t$ be the trace added afterward.
Both are nonnegative.  Since $\lVert b_t\rVert\le1$,
\begin{equation}
 \sum_{t=1}^n g_t\le\sum_{t=1}^n Q_{d+t,d+t},
 \qquad
 \sum_{t=1}^n\ell_t
   =\tr L_0+\sum_{t=1}^n g_t-\tr L_n
   \le\tr Q.
 \label{eq:trace-budgets}
\end{equation}
These two budgets are independent of the length of the sequence.

Fix the following piecewise linear cutoff on $(0,\infty)$:
\[
 f_0(x)=
 \begin{cases}
  0,&0<x\le\tfrac12,\\
  2x-1,&\tfrac12<x<1,\\
  1,&1\le x\le2,\\
  2-\tfrac{x}{2},&2<x<4,\\
  0,&x\ge4.
 \end{cases}
\]
Thus $0\le f_0\le1$ and its Lipschitz constant is $2$.
We use the scalar functions
\begin{equation}
 \begin{gathered}
 a=\tfrac32,\qquad
 w(x)=\min\{x^a,x^{-1}\},\qquad
 p(x)=\int_0^x\frac{w(u)}{u}\,du,\\
 p_1(x)=\int_0^x\frac{u^{a-1}}{(1+u)^a}\,du,
 \qquad
 \phi(x)=\min\{x^{a-1},x^{-1}\}.
 \end{gathered}
 \label{eq:scalar-potentials}
\end{equation}
The potential $p$ will control changes of the cutoff between endpoint
matrices, while $p_1$ will control gains accumulated through intervening
contractions.  The function $\phi$ measures the contribution of each
spectral direction to the size of a scale.  Matrix functions below are
defined by spectral calculus.

For each dyadic scale $\theta\in2^{\mathbb Z}$, set
\[
 X_t=L_t/\theta,
 \qquad
 F_t^\theta=R_t^\top f_0(X_t)R_t,
 \qquad
 k_t(\theta)=\tr\phi(X_t).
\]
Here $R_t$ is $d\times m$, $X_t$ is $d\times d$, and
$F_t^\theta$ is $m\times m$.  By~\eqref{eq:speed-recursion},
$R_t^\top$ is an isometry, so the lift preserves the localized state norm:
\[
 \lVert F_t^\theta x\rVert
 =\lVert f_0(X_t)R_tx\rVert\qquad(x\in\R^m).
\]
It also removes dependence on state coordinates: replacing $R_t$ by
$OR_t$, with $O$ orthogonal, replaces $X_t$ by $OX_tO^\top$ and leaves
$F_t^\theta$ unchanged.
The notation $X_t$ always refers to the scale currently under
consideration; at a fixed scale we also abbreviate $F_t=F_t^\theta$
and $k_t=k_t(\theta)$.
The geometric decay of $\phi$ at zero and infinity gives
\begin{equation}
 \sum_{\theta\in2^{\mathbb Z}}k_t(\theta)\le C d.
 \label{eq:scale-size}
\end{equation}
Indeed, for each positive eigenvalue $\lambda$ of $L_t$, the terms
$\phi(\lambda/\theta)$ on either side of $\theta=\lambda$ form
bounded geometric sums.  On the support band of $f_0$, one has
$\phi\ge1/4$.  It follows that
\begin{equation}
 \begin{gathered}
 \tr f_0(X_t)^2\le4k_t,
 \qquad
 \tr\bigl(X_tf_0(X_t)^2\bigr)\le16k_t,
 \qquad \lVert F_t\rVert_{\op}\le1,\\
 F_t\ne0\quad\Longrightarrow\quad k_t\ge\tfrac14.
 \end{gathered}
 \label{eq:cutoff-moments}
\end{equation}
Every eigenvalue of $L_t$ lies
between the smallest and largest diagonal entries of $Q$.
Consequently only finitely many scales have $F_t^\theta\ne0$ for some
$t$, and every such scale satisfies $\theta\le2\lVert Q\rVert_{\op}$.

We now assign an additive energy to intervals of sequence indices.
Put $h=p+p_1$ and $\widehat X_i=C_iX_{i-1}C_i^\top$.  Define
\begin{equation}
 E_\theta(s,t)=\sum_{i=s+1}^t
 \left\{
   \tr h(X_{i-1})-\tr h(\widehat X_i)
   +\tr h(X_i)-\tr h(\widehat X_i)
 \right\},
 \qquad 0\le s\le t\le n.
 \label{eq:energy-definition}
\end{equation}
The ordered eigenvalues decrease under a contraction and increase
under a positive semidefinite gain.  Since $h$ is increasing, every
trace difference in braces is nonnegative.  In particular,
$E_\theta$ is nonnegative and additive on adjacent index intervals.

The derivatives of the two functions satisfy, uniformly for $x>0$,
\[
 \sum_{\theta\in2^{\mathbb Z}}
   \bigl(p'(x/\theta)+p_1'(x/\theta)\bigr)\le C.
\]
For $p'$, use the bounds $u^{1/2}$ for $u\le1$ and $u^{-2}$ for
$u\ge1$; for $p_1'$, use $u^{1/2}$ and $u^{-1}$ respectively.
All four dyadic sums are geometric.  Pairing ordered eigenvalues
before and after each contraction or gain, and integrating the last
bound between each pair, proves
\begin{equation}
 \sum_{\theta\in2^{\mathbb Z}}\theta E_\theta(0,n)
   \le C\left(\sum_i\ell_i+\sum_i g_i\right)
   \le C\tr Q.
 \label{eq:scale-energy}
\end{equation}
The factor $\theta$ cancels the factor $1/\theta$ from differentiating
$p(x/\theta)$ or $p_1(x/\theta)$.

Here is the variation estimate to be used in the next section.

\begin{lemma}[Weighted increment estimate]
\label{lem:increment}
There are absolute constants $e_0>0$ and $C<\infty$ with the following
property.  For every recursion in Proposition~\ref{prop:filtered},
every positive definite diagonal $Q$, every dyadic $\theta$, and
$0\le s\le t\le n$, if $E=E_\theta(s,t)\le e_0$, then
\begin{equation}
 \bigl\lVert(F_t^\theta-F_s^\theta)
          (I_m+Q/\theta)^{1/2}\bigr\rVert_{\HS}^2
   \le C\bigl(1+k_t(\theta)\bigr)E^{1/a},
 \qquad a=\tfrac32.
 \label{eq:increment}
\end{equation}
\end{lemma}

We prove this in three steps.  A scalar inequality first controls
changes of spectral cutoffs.  Exact matrix expansions then compare
the two endpoints using the old part of the state and the accumulated
new gains.  Finally, a separate estimate bounds those accumulated
gains by the additive energy, even through noncommuting contractions.

\subsection{A scalar comparison}

For $x,y>0$, define
\[
 \mathcal J(x,y)=p(x)-p(y)
      +\left(\frac yx-1\right)\frac{w(x)+w(y)}2.
\]

\begin{lemma}
\label{lem:scalar-comparison}
There is an absolute $c>0$ such that, for all $x,y>0$,
\begin{equation}
 \begin{split}
 \mathcal J(x,y)&\ge c(1+y)|f_0(x)-f_0(y)|^2,\\
 \frac yx\mathcal J(y,x)&\ge c(1+y)|f_0(x)-f_0(y)|^2.
 \end{split}
 \label{eq:scalar-comparison}
\end{equation}
\end{lemma}

\begin{proof}
The function $u\mapsto\log w(e^u)=\min\{au,-u\}$ is
$a$-Lipschitz.  Write $r=y/x$.  For $r<1$, the two endpoint lower
bounds give
\[
 w(u)\ge
 \max\{w(x)(u/x)^a,\,w(y)(y/u)^a\},\qquad y\le u\le x.
\]
The two expressions cross inside $[y,x]$, by the same Lipschitz
bound.  Integrating them against $du/u$ yields
\[
 p(x)-p(y)\ge
 \frac{w(x)+w(y)-2\sqrt{w(x)w(y)}\,r^{a/2}}a.
\]
For $r>1$, integration of the minimum of the corresponding endpoint
upper bounds on $[x,y]$ gives
\[
 p(y)-p(x)\le
 \frac{2\sqrt{w(x)w(y)}\,r^{a/2}-w(x)-w(y)}a.
\]
Using $2\sqrt{w(x)w(y)}\le w(x)+w(y)$ in either case, we obtain
\begin{equation}
 \mathcal J(x,y)\ge\frac{w(x)+w(y)}2 B(r),
 \qquad
 B(r)=r-1-\frac{r^{a/2}-1}{a/2}.
 \label{eq:scalar-remainder}
\end{equation}
If $b=a/2=3/4$, then $B(1)=B'(1)=0$ and
$B''(r)=(1-b)r^{b-2}>0$.  Moreover
$B(0+)=1/b-1>0$ and $B(r)/r\longrightarrow1$ at infinity.
Thus, with an absolute positive constant,
\[
 B(r)\ge c
 \begin{cases}
  1,&0<r\le\tfrac12,\\
  (r-1)^2,&\tfrac12\le r\le2,\\
  r,&r\ge2.
 \end{cases}
\]
The function $rB(1/r)$ satisfies the same bounds, after changing $c$.
Applying \eqref{eq:scalar-remainder} with the arguments reversed
therefore gives these bounds for both left sides of
\eqref{eq:scalar-comparison}.

If $f_0(x)\ne f_0(y)$, at least one argument belongs to
$[1/2,4]$.  When $1/2\le r\le2$, both arguments lie in $[1/4,8]$,
where $w$ is bounded below and $1+y$ is bounded above.  The Lipschitz
bound $|f_0(x)-f_0(y)|\le2x|r-1|$ applies.  When $r\le1/2$,
one $w$ value is bounded below and $y\le4$.  When $r\ge2$,
one $w$ value is again bounded below, and either $y\le4$ or
$y=rx\le4r$.  The preceding three bounds, together with
$|f_0(x)-f_0(y)|\le1$, prove the result.  If the cutoff values
are equal, the right sides vanish and the result follows directly
from \eqref{eq:scalar-remainder}.
\end{proof}

\subsection{Comparing the endpoints}

Fix a scale and indices $s\le t$.  Abbreviate
\[
 Y=X_s,\qquad X=X_t,\qquad
 \mathcal C=C_t\cdots C_{s+1},\qquad Z=\mathcal C Y\mathcal C^\top,
 \qquad V=X-Z.
\]
An empty product is the identity.  The recursion shows that $V$ is
positive semidefinite: it is the sum of all intervening gains,
each transported through the contractions that follow it.
The matrices $Y,X,Z$ are positive definite, since every $C_i$ is
invertible.

\begin{lemma}[Endpoint comparison]
\label{lem:endpoint-comparison}
With the preceding notation,
\begin{equation}
 \bigl\lVert(F_t-F_s)(I_m+Q/\theta)^{1/2}\bigr\rVert_{\HS}^2
 \le C\left\{
   \tr\bigl(p(Y)-p(Z)\bigr)
   +\tr\bigl(p(X)-p(Z)\bigr)
 \right\}.
 \label{eq:endpoint-comparison}
\end{equation}
\end{lemma}

\begin{proof}
Let $P_s$ be the coordinate projection onto the initial block and
the first $s$ new coordinates, and put
\[
 Q_{\rm old}=P_s(Q/\theta)P_s,
 \qquad Q_{\rm new}=Q/\theta-Q_{\rm old}.
\]
Diagonality of $Q$ makes these nonnegative matrices with disjoint
coordinate supports.  Since $R_tP_s=\mathcal C R_s$, the support identities are
\begin{equation}
 \begin{gathered}
 R_tR_s^\top=\mathcal C,\qquad
 R_tQ_{\rm old}R_s^\top=\mathcal C Y,\qquad
 R_tQ_{\rm old}R_t^\top=Z,\\
 R_tQ_{\rm new}R_t^\top=V,\qquad F_sQ_{\rm new}=0.
 \end{gathered}
 \label{eq:old-new-support}
\end{equation}
We split the right weight into $I_m+Q_{\rm old}$ and
$Q_{\rm new}$, and use the intermediate matrix
$\widetilde F=R_t^\top f_0(Z)R_t$ for the first part.

\paragraph{The old state under contraction.}
Choose orthonormal eigenbases of $Z$ and $Y$, with eigenvalues
$z_i$ and $y_j$.  Let $m_{ij}$ be the squared entries of $\mathcal C$ in
these bases.  Its row and column sums are at most one; denote their
deficits by
\[
 d_i=1-\sum_jm_{ij},\qquad h_j=1-\sum_im_{ij}.
\]
The matrix $n_{ij}=m_{ij}y_j/z_i$ has both row and column sums one.
For rows this follows from $Z=\mathcal C Y\mathcal C^\top$; for columns it follows
from the identity
\[
 \mathcal C^\top Z^{-1}\mathcal C=Y^{-1},
\]
which uses invertibility of $\mathcal C$.

Expanding the weighted squared distance with
\eqref{eq:old-new-support}, its two diagonal terms are
$\tr(f_0(Z)^2(I_d+Z))$ and $\tr(f_0(Y)^2(I_d+Y))$.
The cross term, before multiplication by two, is
$\sum_{ij}m_{ij}(1+y_j)f_0(z_i)f_0(y_j)$.  Since
$\sum_jm_{ij}y_j=z_i$, this gives the exact expansion
\begin{equation}
 \begin{split}
 &\bigl\lVert(\widetilde F-F_s)
                  (I_m+Q_{\rm old})^{1/2}\bigr\rVert_{\HS}^2\\
 &\quad=\sum_{ij}m_{ij}(1+y_j)
                   \bigl(f_0(z_i)-f_0(y_j)\bigr)^2
       +\sum_i d_i f_0(z_i)^2
       +\sum_j h_j(1+y_j)f_0(y_j)^2.
 \end{split}
 \label{eq:contraction-distance}
\end{equation}
Expanding $\mathcal J$ instead, the two stochastic sums of $n$
and the relation $n_{ij}z_i/y_j=m_{ij}$ yield
\begin{equation}
 \tr\bigl(p(Y)-p(Z)\bigr)
  =\sum_{ij}n_{ij}\mathcal J(y_j,z_i)
    +\frac12\sum_i d_i w(z_i)
    +\frac12\sum_j h_j w(y_j).
 \label{eq:contraction-deficits}
\end{equation}
The second inequality of Lemma~\ref{lem:scalar-comparison} controls
the double sum in \eqref{eq:contraction-distance}.  The remaining
terms are controlled by the deficits in
\eqref{eq:contraction-deficits}, since both $f_0(u)^2$ and
$(1+u)f_0(u)^2$ are bounded by an absolute multiple of $w(u)$.
This proves the required bound for $\widetilde F-F_s$.

\paragraph{The accumulated gain.}
Now choose eigenbases of $X$ and $Z$, with eigenvalues $x_i$ and
$z_j$, and let $m_{ij}$ be the squared entries of their orthogonal
change of basis.  Both sums of $m$ equal one.  Set
$n_{ij}=m_{ij}z_j/x_i$.  Its row sums are at most one by
$Z\preceq X$, and its column sums are at most one by
$X^{-1}\preceq Z^{-1}$: explicitly, for a unit $z_j$-eigenvector
$v_j$,
\[
 \sum_i n_{ij}=z_jv_j^\top X^{-1}v_j
       \le z_jv_j^\top Z^{-1}v_j=1.
\]
Write $d'_i,h'_j$ for these row and column deficits.  Expansion of
$\mathcal J$ now gives
\begin{equation}
 \tr\bigl(p(X)-p(Z)\bigr)
  =\sum_{ij}m_{ij}\mathcal J(x_i,z_j)
    +\frac12\sum_i d'_i w(x_i)
    +\frac12\sum_j h'_j w(z_j).
 \label{eq:gain-deficits}
\end{equation}
The coisometry $R_tR_t^\top=I_d$ and
\eqref{eq:old-new-support} show that
\begin{align*}
 &\bigl\lVert(F_t-\widetilde F)
                   (I_m+Q_{\rm old})^{1/2}\bigr\rVert_{\HS}^2\\
 &\qquad=\bigl\lVert(f_0(X)-f_0(Z))(I_d+Z)^{1/2}
                           \bigr\rVert_{\HS}^2
   =\sum_{ij}m_{ij}(1+z_j)\bigl(f_0(x_i)-f_0(z_j)\bigr)^2.
\end{align*}
The first inequality of Lemma~\ref{lem:scalar-comparison} controls
this by \eqref{eq:gain-deficits}.

It remains to include the new-coordinate part of the right weight.
By \eqref{eq:old-new-support}, it is exactly
\begin{equation}
 \bigl\lVert(F_t-F_s)Q_{\rm new}^{1/2}\bigr\rVert_{\HS}^2
   =\tr\bigl(f_0(X)^2V\bigr)
   =\sum_i x_i d'_i f_0(x_i)^2.
 \label{eq:new-coordinate-distance}
\end{equation}
Here $x_id'_i$ is the $i$th diagonal entry of $X-Z$ in the
$X$-eigenbasis.  The bound $u f_0(u)^2\le Cw(u)$ controls this
term by the row deficits in \eqref{eq:gain-deficits}.
Finally apply the squared triangle inequality to
$(F_t-\widetilde F)+(\widetilde F-F_s)$ with the old weight,
and add \eqref{eq:new-coordinate-distance}.  This proves
\eqref{eq:endpoint-comparison}.
\end{proof}

The endpoint comparison has accounted for all coordinate weights.
Its remaining difficulty is that $p(X)-p(Z)$ describes a single
accumulated gain, whereas the energy in
\eqref{eq:energy-definition} records gains separated by contractions.
The next estimate connects these two quantities.

\subsection{Controlling the accumulated gain}

\begin{lemma}
\label{lem:direct-gain}
For the matrices $X,Y,Z,V$ above, suppose
$E=E_\theta(s,t)\le e_0$, where
\[
 c_a=a2^a,\qquad e_0=\frac1{16c_a},\qquad a=\tfrac32.
\]
If $q=(c_aE)^{1/a}$, then $q<1/4$ and
\begin{equation}
 V\preceq q(I_d+X),
 \qquad
 \tr\bigl(p(X)-p(Z)\bigr)\le Cq\tr\phi(X).
 \label{eq:direct-gain}
\end{equation}
\end{lemma}

\begin{proof}
Follow the old and new parts separately through the interval
$[s,t]$.  Start with $(Z_*,V_*)=(Y,0)$.  At each contraction,
replace both matrices by their congruences under $C_i$.  At each
gain, add the positive semidefinite gain to $V_*$ only, interpolating
linearly during this addition.  At every stage
$D=Z_*+V_*$ is the matrix in the original contraction-and-gain path;
at the endpoint $(Z_*,V_*)=(Z,V)$.

Define
\[
 S=I_d+Z_*,\qquad A=S^{-1/2}V_*S^{-1/2},\qquad
 G=\tr A^a.
\]
The use of $I_d+Z_*$ makes this quantity stable under contractions.
We will show that its increases during gains are paid for by the
$p_1$ part of the energy.

\paragraph{Contractions do not increase $G$.}
For a contraction $C$, put $S'=I_d+CZ_*C^\top$ and
$T_c=(S')^{-1/2}CS^{1/2}$.  The new relative matrix is
$A'=T_cAT_c^\top$.  Moreover
\[
 S'-CSC^\top=I_d-CC^\top\succeq0,
 \qquad T_cT_c^\top\preceq I_d.
\]
Thus the ordered eigenvalues of $A'$ are no larger than those of
$A$; for example, compare the singular values of $T_cA^{1/2}$
and $A^{1/2}$.  Raising these nonnegative eigenvalues to the power
$a$ and summing proves $\tr(A')^a\le\tr A^a$.
This argument uses eigenvalue comparison, not operator monotonicity
of the power $a$.

\paragraph{The derivative during a gain.}
Here $S$ is fixed and $\dot V_*\succeq0$.  Introduce
\[
 J=I_d+D=S+V_*,\qquad
 H=J^{-1/2}V_*J^{-1/2},\qquad
 T=S^{-1/2}J^{1/2}.
\]
Since $I_d-H=J^{-1/2}SJ^{-1/2}$ is positive definite,
\[
 T^\top T=(I_d-H)^{-1}.
\]
Write the polar decomposition as
$T=U(I_d-H)^{-1/2}$, where $U$ is orthogonal.  Then
\begin{equation}
 A=THT^\top
   =U\bigl[H(I_d-H)^{-1}\bigr]U^\top.
 \label{eq:relative-spectrum}
\end{equation}
In particular, the eigenvalues of $A$ are $h/(1-h)$, where $h$
ranges over the eigenvalues of $H$.

For a continuously differentiable scalar function $g$ on the spectral
range of a differentiable symmetric matrix path $A(\xi)$,
$\frac{d}{d\xi}\tr g(A)=\tr(g'(A)\dot A)$. One can verify this first
for polynomials by cyclicity of the trace, then approximate $g'$ uniformly
by polynomials and integrate along the path. Applied to $g(u)=u^a$, this gives
\begin{equation}
 \dot G
   =a\tr\bigl(S^{-1/2}A^{a-1}S^{-1/2}\dot V_*\bigr).
 \label{eq:gain-derivative}
\end{equation}
This formula remains valid when $A$ is singular, because
$u\mapsto u^a$ is continuously differentiable at zero for $a>1$.
The coefficient in this trace has the exact conjugation identity
\begin{align}
 J^{1/2}S^{-1/2}A^{a-1}S^{-1/2}J^{1/2}
   &=T^\top A^{a-1}T\notag\\
   &=H^{a-1}(I_d-H)^{-a}.
 \label{eq:gain-conjugation}
\end{align}
The second equality follows from \eqref{eq:relative-spectrum} and
the polar decomposition.  The factors on its right commute because
they are functions of $H$; no commutation of $Z_*$ and $V_*$ is
assumed.

As long as $G\le1$, every eigenvalue of $A$ is at most one, so
\eqref{eq:relative-spectrum} gives $H\preceq I_d/2$.  Since
$a-1=1/2$, spectral calculus in $H$ then gives
\[
 H^{a-1}(I_d-H)^{-a}\preceq2^aH^{1/2}.
\]
Also $V_*\preceq D$, and hence
\[
 H\preceq J^{-1/2}DJ^{-1/2}=D(I_d+D)^{-1}.
\]
The square root preserves positive semidefinite order.  One direct
proof is the convergent integral formula
\[
 M^{1/2}=\frac1\pi\int_0^\infty
       u^{-1/2}M(uI_d+M)^{-1}\,du,
       \qquad M\succeq0:
\]
the integrand preserves order because
$M(uI_d+M)^{-1}=I_d-u(uI_d+M)^{-1}$ and inversion reverses order.
It follows that
\[
 H^{1/2}\preceq D^{1/2}(I_d+D)^{-1/2}.
\]
Conjugating \eqref{eq:gain-conjugation} back by $J^{-1/2}$ yields
\[
 S^{-1/2}A^{a-1}S^{-1/2}
   \preceq2^aD^{1/2}(I_d+D)^{-3/2}=2^ap_1'(D).
\]
Pair this inequality with $\dot V_*\succeq0$ in
\eqref{eq:gain-derivative}.  Since $\dot D=\dot V_*$ during a gain,
we obtain the conditional differential bound
\begin{equation}
 \dot G\le c_a\tr\bigl(p_1'(D)\dot V_*\bigr)
       =c_a\frac{d}{d\xi}\tr p_1(D),
       \qquad G\le1.
 \label{eq:gain-differential-bound}
\end{equation}

\paragraph{Closing the bound.}
Initially $G=0$, and contractions cannot increase it.  The sum of
all increases of $\tr p_1(D)$ during gains is at most $E$, by
\eqref{eq:energy-definition}.  If $G$ first reached one during a
gain, integration of \eqref{eq:gain-differential-bound} up to that
point would give $G\le c_aE\le1/16$, a contradiction.
Thus the differential bound applies throughout, and
$G\le c_aE$ at the endpoint.  In particular,
\[
 A\preceq qI_d,\qquad
 V\preceq q(I_d+Z)\preceq q(I_d+X),\qquad
 q\le16^{-2/3}<\tfrac14.
\]
If $E=0$, this already gives $V=0$ and the second conclusion in
\eqref{eq:direct-gain}.  Suppose now that $q>0$.

The ordered eigenvalues of $X$ and $Z$ obey
\[
 x_j\ge z_j\ge(1-q)x_j-q,
\]
because $X-Z=V\preceq q(I_d+X)$.  For a corresponding pair $x,z$,
if $x\le4q\le1$, then
\[
 p(x)-p(z)\le p(x)=\frac{x^a}{a}
       \le\frac4a qx^{a-1}=\frac4a q\phi(x).
\]
If $x>4q$, then $z\ge x/2$ and $x-z\le q(1+x)$.
On $[x/2,x]$ we have $w(u)\le2^aw(x)$ and $u^{-1}\le2/x$,
so
\[
 p(x)-p(z)\le Cq(1+x^{-1})w(x)\le Cq\phi(x).
\]
The last inequality holds for both $x\le1$ and $x\ge1$ directly
from \eqref{eq:scalar-potentials}.  Summing over the eigenvalues
proves the trace bound in \eqref{eq:direct-gain}.
\end{proof}

\begin{proof}[Proof of Lemma~\ref{lem:increment}]
The total variation of $\tr p$ along the separate contraction and
gain steps from $s$ to $t$ is at most $E$.  Therefore
\[
 \bigl|\tr p(Y)-\tr p(X)\bigr|\le E.
\]
Together with Lemma~\ref{lem:direct-gain}, this bounds the sum on
the right of \eqref{eq:endpoint-comparison} by
\[
 E+2\tr\bigl(p(X)-p(Z)\bigr)
       \le E+CE^{1/a}k_t.
\]
Since $0\le E\le e_0\le1$, we have $E\le E^{1/a}$.
Lemma~\ref{lem:endpoint-comparison} now proves
\eqref{eq:increment}, including $E=0$.
\end{proof}

The two estimates needed for the probability argument are now in
place.  Equation~\eqref{eq:scale-energy} bounds the sum of the
energies over all spectral scales, while
Lemma~\ref{lem:increment} turns a short interval in this energy into
a small Gaussian increment.  In particular, zero energy implies
$F_t^\theta=F_s^\theta$ exactly, so repeated indices at zero
distance will not increase the size of a covering in the next section.
